\documentclass[12pt]{article}
\usepackage[left = 0.6in, right = 0.6in, top = 0.6in, bottom = 1.5in]{geometry} 
\usepackage{enumerate, pgfplots, fontawesome5, booktabs, xcolor, hyperref, amsmath,amsthm,amssymb,amsfonts, fancyhdr, color, comment, graphicx, environ, actuarialsymbol}
\pagestyle{fancy}
\setlength{\headheight}{65pt}
%%%%%%%%%%%%%%%%%%%%%%%%%%%%%%%%%%%%%%%%%%%%%
\lhead{}
\rhead{functions 2}
%%%%%%%%%%%%%%%%%%%%%%%%%%%%%%%%%%%%%%%%%%%%%
%%%%%%%%%%%%%%%%%%%%%%%%%%%%%%%%%%%%%%%%%%%%%
\newcommand{\emptygrid}[3]{ % #1: range, #2: y-label, #3: "labels" or "nolabels"
    \begin{tikzpicture}
    \begin{axis}[
        axis lines = middle,
        xlabel = $x$,
        ylabel = $#2$,
        xmin=-#1, xmax=#1,
        ymin=-#1, ymax=#1,
        grid = both,
        grid style = {line width=.1pt, draw=gray!30},
        xtick distance = 5,
        ytick distance = 5,
        minor tick num = 4,
        tick label style = {font=\tiny},
        width=8cm, height=8cm,
        axis equal image,
        % Logic to hide labels if #3 is "nolabels"
        xticklabels={ \ifx\empty#3\empty\else\ifnum\pdfstrcmp{#3}{nolabels}=0 \else \fi\fi },
        yticklabels={ \ifx\empty#3\empty\else\ifnum\pdfstrcmp{#3}{nolabels}=0 \else \fi\fi }
    ]
    \end{axis}
    \end{tikzpicture}
}
%%%%%%%%%

\begin{document}
\section{Inverse functions}
An inverse of a function $f$ is another function $f^{-1}$ that does the opposite operations in reverse order
\begin{itemize}
    \item When can inverses exist?
    \vspace{3cm}

    \item $f^{-1}(f(x)) = f(f^{-1}(x)) = x$
    \item How do you find inverse? \textbf{ALWAYS} follow \textit{this} process if you want full marks:
    \begin{enumerate}
        \item "Let $y = f^{-1}(x)$"
        \item "Then $x = f(y)$" 
        \item (solve for $y$)
        \item "Therefore $y = f^{-1}(x) = \dots$"
    \end{enumerate}

    \begin{enumerate}[(i)]
        \item Let $f: D \to \mathbb{R}, f(x) = x^{2}+6x - 3$
        \begin{enumerate}
            \item Find the maximal domain $D$ such that $f$ has an inverse
            \item Find the inverse $f^{-1}$
            \item Find the slope of the tangent to $f^{-1}$ at $x = 10$
        \end{enumerate}
        \vspace{5cm}

        \item The function defined by $f : A \to \mathbb{R}$, where $f(x) = -2(x+3)^2 + 4$ will have an inverse function if its domain is
        \begin{enumerate}
            \item[A.] $\mathbb{R}$
            \item[B.] $\mathbb{R}^- \cup \{0\}$
            \item[C.] $\{x : x < -4\}$
            \item[D.] $\{x : x \geq -2\}$
            \item[E.] $\{x : x \leq 3\}$
        \end{enumerate}
    \end{enumerate}
    
    \item POI between $f$ and $f^{-1}$? Can lie on 
    \begin{align*}
        y &= x \\
        y &= -x + c, \quad \text{for some }c
    \end{align*}
    \item What do the POI on each of those look like? Either $(a,a)$ [on $y = x$] or in pairs of $(a,b),(b,a)$ [on $y = -x + c$]
    
    \item $f(x) = x$ gives solutions to all $f(x) = f^{-1}(x)$ ONLY IF $f$ is strictly increasing

\end{itemize}


    \begin{enumerate}
    \item (see 2003 VCE Mathematical Methods Exam 2 Part 2 Question 4)

    \begin{enumerate}
        \item[a.] Sketch the graph of the equation $x = \frac{3t}{5+t^2}, \ t \geq 0$.
        \item[b.] 
        \begin{enumerate}
            \item[i.] Find, correct to 4 decimal places, the least value of $a$ such that the function 
            \[ h: [a, \infty) \to \mathbb{R}, \ h(t) = \frac{3t}{5+t^2} \]
            has an inverse function.
            \item[ii.] For this value of $a$, sketch the graph of $h^{-1}$. Label any endpoint with its coordinates. Label any asymptote with its equation.
        \end{enumerate}
        \item[c.] Find the rule for $h^{-1}$.
    \end{enumerate}

    \item Consider the function $f: [2, \infty) \to \mathbb{R}, \ f(x) = \sqrt{x^2 - p}$.

        \begin{enumerate}
            \item[a.] Find the value of $p$ if $f$ has an inverse function $f^{-1}: [1, \infty) \to \mathbb{R}$.
            \item[b.] Find the rule for $f^{-1}$.
        \end{enumerate}

    \item Determine the value(s) of $k$ such that the inverse function of $f: [-k, \infty) \to \mathbb{R}, \ f(x) = kx^2 + x + 1$ is given by the rule $f^{-1}(x) = \sqrt{2x - 1} - 1$.

    \item Let $f: [a, 2) \to \mathbb{R}, \ f(x) = x^3 - 2x^2 - 4x + 5$ where $a$ is the smallest value for which the inverse function $f^{-1}$ exists. Find the exact value of $a$.
\end{enumerate}

\end{document}